\scanpage{37}
\mspage{17}
Comment remonter les \inserted{2-}sous-groupes \inserted{de Sylow} $\simeq\mathbb Z/2\times\mathbb Z/2$ de $\Gamma$, engendrés resp. par $\{\sigma_0,\tau\},\{\sigma_1,\tau\},\{\sigma_\infty,\tau\}$? Et bien justement, ils ne se remontent \uline{pas} --- on sait bien que \unsure{les extensions locales|l'extension locale} de $\Gamma_0^+$ \inserted{(\unsure{précis})} par $\mathbb Z_0$ (quotient de $\mathbb Z^{\{0,1,\infty\}}$ par $\mathbb Z^{\{1,\infty\}}$) n'est \uline{pas} scindée, i.e. le générateur $\dot\sigma_0$ de $\Gamma_0^+$ ne se remonte \uline{pas} en un élément d'ordre $2$ de $\mathcal E$ qui centralise (ou ce qui revient au $=$, normalise) $L_0$.

On pourrait s'amuser à écrire $S\Gamma$ et $S\mathcal E$ en générateurs et relations --- je ne vais pas m'y attarder. Par contre, je vais donner des formules pour l'action des $\operatorname{int}(\tilde\rho)$ et des $\operatorname{int}(\tilde\sigma_i)$. On trouve
\[
(15)\qquad(\operatorname{int}\tilde\rho)(p,\sigma,g_0,g_1,g_\infty)=(p,\dot\rho\sigma\dot\rho^{-1},\rho(g_\infty),\rho(g_0),\rho(g_1))
\]
\[
(16)\qquad\operatorname{int}(\tilde\sigma_0)(p,\sigma,g_0,g_1,g_\infty)=(p,\dot\sigma_0\sigma\dot\sigma_0^{-1},\sigma_0(g_0),\sigma_0(g_\infty),\sigma_0(g_1))
\]
On voit donc: a) Pour que $\tilde u=(p,\unsure{\sigma},g_0,g_1,g_\infty)$ commute à $\rho$, il f. et s. que \deleted{$g_0,g_1,g_\infty$ \unsure{soi}} $\sigma$ commute à $\rho$ i.e. $\sigma\in\{1,\rho,\rho^2\}$, \uline{et} qu'on ait
\[
(17)\qquad g_1=\rho(g_0),\quad g_\infty=\rho(g_1),\quad g_0=\rho(g_\infty)
\]
Ce qui signifie aussi que $g_1,g_\infty$ se déduisent de $g_0$ par $g_1=\rho(g_0),\ g_\infty=\rho(g_1)=\rho^2(g_0)=\rho^{-1}(\unsure{g_0|g_\infty})$
\mspage{36}
